\Needspace{15\baselineskip}
\Article{25}

\textsc{Theorem.}\quad \textit{Let $n$ be the product of two coprime positive integers $a$, $b$,} \textit{and put}

\[
\begin{aligned}
P&=1+ r^{aa}+r^{4aa}+r^{9aa}+\text{etc.}+r^{(b-1)^2aa}\\
Q&=1+r^{bb}+r^{4bb}+r^{9bb}+\text{etc.}+r^{(a-1)^2bb}
\end{aligned}
\]

\textit{Then I claim that $W=PQ$.}

\textit{Proof.}\quad Let $\alpha$ denote indefinitely the numbers $0,1,2,3\dotsfour a-1$, $\beta$ the numbers $0,1,2,3\dotsfour b-1$, and $\nu$ the numbers $0,1,2,3\dotsfour n-1$.\quad Then evidently

\[
\begin{aligned}
P=\Sigma r^{aa\beta\beta},\quad Q=\Sigma r^{bb\alpha\alpha},\quad W=\Sigma r^{\nu\nu}
\end{aligned}
\]

Thus $PQ=\Sigma r^{aa\beta\beta+bb\alpha\alpha}$, on substituting all the values of $\alpha$ and $\beta$, combined with one another in every way. Further, since $2ab\alpha\beta=2\alpha\beta n$, we have $PQ=\Sigma r^{(a\beta+b\alpha)^2}$. But it is easy to see that the individual values of $a\beta+b\alpha$ are different from one another and equal to some value of $\nu$.\quad Hence $PQ=\Sigma r^{\nu\nu}=W$.

It should also be noted that $r^{aa}$ is a proper root of the equation $x^b-1=0$, and $r^{bb}$ a proper root of the equation $x^a-1=0$.

\SourceLeaf{47}

\Needspace{10\baselineskip}
\Article{26}

Now let $n$ be the product of three pairwise coprime numbers $a,b,c$. It is clear that, on putting $bc=b'$, $a$ and $b'$ will also be coprime; hence $W$ is the product of the two factors

\[
\begin{aligned}
&1+r^{aa}+r^{4aa}+r^{9aa}+\text{etc.}+r^{(b'-1)^2aa}\\
&1+r^{b'b'}+r^{4b'b'}+r^{9b'b'}+\text{etc.}+r^{(a-1)^2b'b'}
\end{aligned}
\]

But since $r^{aa}$ is a proper root of the equation $x^{bc}-1=0$, the first factor itself will be the product of

\[
\begin{aligned}
&1+\rho^{bb}+\rho^{4bb}+\rho^{9bb}+\text{etc.}+\rho^{(c-1)^2bb}\\
&1+\rho^{cc}+\rho^{4cc}+\rho^{9cc}+\text{etc.}+\rho^{(b-1)^2cc}
\end{aligned}
\]

if we put $r^{aa}=\rho$.\quad It is therefore clear that $W$ is the product of the three factors

\[
\begin{aligned}
&1+r^{bbcc}+r^{4bbcc}+r^{9bbcc}+\text{etc.}+r^{(a-1)^2bbcc}\\
&1+r^{aacc}+r^{4aacc}+r^{9aacc}+\text{etc.}+r^{(b-1)^2aacc}\\
&1+r^{aabb}+r^{4aabb}+r^{9aabb}+\text{etc.}+r^{(c-1)^2aabb}
\end{aligned}
\]

where $r^{bbcc}$, $r^{aacc}$, $r^{aabb}$ are respectively proper roots of the equations $x^a-1=0$, $x^b-1=0$, $x^c-1=0$.

\Needspace{12\baselineskip}
\Article{27}

It is easy to conclude in general that, if $n$ is the product of any number of pairwise coprime factors $a,b,c$, etc., $W$ is the product of the same number of factors, namely

\[
\begin{aligned}
&1+r^{\frac{nn}{aa}}+r^{\frac{4nn}{aa}}+r^{\frac{9nn}{aa}}+\text{etc.}+r^{\frac{(a-1)^2nn}{aa}}\\
&1+r^{\frac{nn}{bb}}+r^{\frac{4nn}{bb}}+r^{\frac{9nn}{bb}}+\text{etc.}+r^{\frac{(b-1)^2nn}{bb}}\\
&1+r^{\frac{nn}{cc}}+r^{\frac{4nn}{cc}}+r^{\frac{9nn}{cc}}+\text{etc.}+r^{\frac{(c-1)^2nn}{cc}},\;\text{etc.}
\end{aligned}
\]

where $r^{\frac{nn}{aa}}$, $r^{\frac{nn}{bb}}$, $r^{\frac{nn}{cc}}$, etc. are proper roots of the equations $x^a-1=0$, $x^b-1=0$. $x^c-1=0$, etc.

\Needspace{14\baselineskip}
\Article{28}

These principles now make the transition to a complete determination of $W$ for any value of $n$ immediately clear.\quad Specifically, resolve $n$ into factors
\SourceLeaf{48}
$a,b,c$, etc., that are either distinct prime numbers or powers of distinct prime numbers, and put $r^{\frac{nn}{aa}}=A$, $r^{\frac{nn}{bb}}=B$, $r^{\frac{nn}{cc}}=C$, etc. Then $A,B,C$, etc. will be proper roots of the equations $x^a-1=0$, $x^b-1=0$, $x^c-1=0$, etc., and $W$ the product of the factors

\[
\begin{aligned}
&1+A+A^4+A^9+\text{etc.}+A^{(a-1)^2}\\
&1+B+B^4+B^9+\text{etc.}+B^{(b-1)^2}\\
&1+C+C^4+C^9+\text{etc.}+C^{(c-1)^2}
\end{aligned}
\]

Each of these factors can be determined from what we showed in Articles $20,21,23$, so the value of the product will also become known.\quad It will be useful to collect the rules for determining these factors in a single overview here.\quad Since the root $A$ is $=\frac{kn}{a}\pdot\frac{360^0}{a}$, the sum $1+A+A^4+A^9+\text{etc.}+A^{(a-1)^2}$, which we shall denote by $L$, will be determined by the number $\frac{kn}{a}$ in the same way as $W$ is determined by $k$ in our general investigation.\quad Twelve cases must now be distinguished.

I.\quad If $a$ is a prime number of the form $4\mu+1$, say $=p$, or a power of such a prime number with an odd exponent, and $\frac{kn}{a}$ is a quadratic residue of $p$, then $L=+\surd a$.

II.\quad If the other conditions remain unchanged and $\frac{kn}{a}$ is a quadratic nonresidue of $p$, then $L=-\surd a$.

III.\quad If $a$ is a prime number of the form $4\mu+3$, say $=p$, or a power of such a prime number with an odd exponent, and $\frac{kn}{a}$ is a quadratic residue of $p$, then $L=+i\surd a$.

IV.\quad If, with the other conditions as in III, $\frac{kn}{a}$ is a quadratic nonresidue of $p$, then $L=-i\surd a$.

V.\quad If $a$ is the square, or a higher power with an even exponent, of an (odd) prime number, then $L=+\surd a$.

VI.\quad If $a=2$, then $L=0$.

VII.\quad If $a=4$, or it is a higher power of two with an even exponent, and $\frac{kn}{a}$ is of the form $4\mu+1$, then $L=(1+i)\surd a$.

VIII.\quad If, with the other conditions as in VII, $\frac{kn}{a}$ is of the form $4\mu+3$, then $L=(1-i)\surd a$.

IX.\quad If $a=8$, or it is a higher power of two with an odd exponent, and $\frac{kn}{a}$ is of the form $8\mu+1$, then $L=(1+i)\surd a$.

\SourceLeaf{49}

X.\quad If, with the other conditions as in IX, $\frac{kn}{a}$ is of the form $8\mu+3$, then $L=(-1+i)\surd a$.

XI.\quad If the other conditions remain unchanged and $\frac{kn}{a}$ is of the form $8\mu+5$, then $L=(-1-i)\surd a$.

XII.\quad If the other conditions remain unchanged and $\frac{kn}{a}$ is of the form $8\mu+7$, then $L=(1-i)\surd a$.

\Needspace{17\baselineskip}
\Article{29}

For example, let $n=2520=8\pdot9\pdot5\pdot7$ and $k=13$.\quad Here we obtain

{\small
\[
\begin{aligned}
&\text{for }a=8,\;\text{by case XII,}\quad L=(1-i)\surd8\\
&\text{for the factor }9,\;\text{by case V, the corresponding sum is }=\surd9\\
&\text{for the factor }5,\;\text{by case II, the corresponding sum is }=-\surd5\\
&\text{for the factor }7,\;\text{by case III, the corresponding sum is }=+i\surd7\mathord{.}
\end{aligned}
\]
}

Thus $W=(1-i)\pdot(-i)\pdot\surd2520=(-1-i)\surd2520$.

For the same value of $n$, let $k=1$; then the corresponding sums are

\[
\begin{aligned}
&\text{for the factor }8\;\text{the sum }(-1+i)\surd8\\
&\text{for the factor }9\;\text{the sum }\surd9\\
&\text{for the factor }5\;\text{the sum }\surd5\\
&\text{for the factor }7\;\text{the sum }-i\surd7
\end{aligned}
\]

This gives the product $W=(1+i)\surd2520$.
